\begin{longtblr}{colspec={Q[l]X[l]}, width=\linewidth, row{even}={bg=rowshade}, rowhead=1}
\toprule
\fort{&nlte_params} & Meaning\\
\midrule
\fort{convlim} & Convergence limit of the relative population change. A positive value requests NLTE iterations; the default 0 keeps the atom in LTE.\\
\fort{convmax} & Convergence limit of the relative population change within one depth patch (default \fort{convlim}/10; 0 is never met). An atom that changes by less than \fort{convmax} in a patch is no longer updated there; with \fort{premature_finish} the patch then stops iterating, and once every patch has stopped the run ends. The closing message then reports that every patch met \fort{convmax} and the iteration count. Keep \fort{convmax} below \fort{convlim}, so that \fort{convlim} decides.\\
\fort{stall_iter} & Number of iterations after which NLTE iterations that are going nowhere stop (default 10; 0 switches the test off). Once the largest relative change of an iteration has not fallen below its smallest earlier value for \fort{stall_iter} iterations, as in a run that oscillates or sits on a floor above \fort{convlim}, the iterations end as after a stop request (Sec.~\ref{sec:stopping}): the departures are written, the output pass runs and the closing warning reports the stall with the smallest change and its iteration. A run that converges, however slowly, keeps reaching new minima and is not stopped. An iteration whose largest change is above \fort{max_update}$-1$, which only a correction cut by the limiter produces, starts the count again, since such a transient can still recover; a run held at the limiter runs on to \fort{maxiter}.\\
\fort{premature_finish} & Freeze the columns whose atoms and all deeper columns have met \fort{convmax} (default \fort{T}).\\
\fort{preconditioning} & Preconditioned rates of Rybicki \& Hummer (1992) (default \fort{T}); \fort{rate_formulation='largest'} or \fort{'original'}.\\
\fort{scattering_preconditioning} & Local approximate operator for continuum scattering (default \fort{F}), limited by \fort{scat_op_max}.\\
\fort{decouple_continuum} & Use separate frequency sets for bound-bound and bound-free radiation (default \fort{T}).\\
\fort{initial_Jnu}, \fort{fixed_Jnu} & Initial radiation field (\fort{'Planck'}) or a fixed value.\\
{\fort{lse_solver},\\\fort{sor_w},\\\fort{cswitch}} & Statistical-equilibrium solver and its damping.\\
{\fort{conv_pop_floor},\\\fort{max_update},\\\fort{min_pop_frac}} & Safeguards of the population update. \fort{max_update} (default 100) is the largest factor by which a level population may change in one iteration, in either direction; a correction cut by it reports a relative change of \fort{max_update}$-1$ or $1-1/$\fort{max_update}, so levels held at the limiter produce ``global'' changes of thousands of percent that cannot fall below \fort{convlim}. \fort{min_pop_frac} (default $10^{-30}$) keeps a level's population above that fraction of the atom's total; for a level whose LTE population is itself below that fraction the floor is $b=1$, since a larger departure coefficient does not fit the single-precision transfer. \fort{conv_pop_floor} (default 0, off) is the population fraction above which a level has full weight in the convergence test; below it the change is scaled down in proportion. It keeps levels that hold almost none of the atom, such as the uppermost levels of the last ion, from deciding \fort{convlim} and \fort{convmax}; check the emergent spectrum when using it. The unweighted change of every level is still written to \File{convergence\_report\_NNN.txt}.\\
\fort{ng_step} & Ng acceleration: $-1$ off, 0 automatic gate (default), $>0$ fixed schedule starting at \fort{ng_start} (default 8).\\
{\fort{maser_chi_floor},\\\fort{suppress_masers}, \dots} & Guards against inverted populations.\\
\bottomrule
\end{longtblr}
